\documentclass[10pt,a4paper]{article}
\usepackage[utf8]{inputenc}
\usepackage[T1]{fontenc}
\usepackage[margin=18mm]{geometry}
\usepackage{longtable}
\usepackage{hyperref}
\usepackage{enumitem}
\setlist[itemize]{noitemsep, topsep=2pt}
\setlength{\parindent}{0pt}
\begin{document}

\section*{Non-Verbal \& Multimodal Discourse Analysis Worksheet}

\textit{Qualitative analysis of pauses, trailing sentences, prosody/intonation, vocalisations (coughs, sighs, scoffs), and kinesics (hand gestures, facial expressions, posture) in political news discourse}\par\smallskip

\bigskip\hrule\bigskip

\subsection*{Guidelines: Five non-verbal \& paralinguistic features to look for in every event}

\begin{longtable}{|p{0.23\textwidth}|p{0.23\textwidth}|p{0.23\textwidth}|p{0.23\textwidth}|}
\hline
Feature & What to look for & Illustrative example (from sample) & Questions to answer \\ \hline \endhead
\textbf{1. Pauses, Silences \& Abrupt Stops} & Micro-pauses \texttt{(.)}, timed silences \texttt{(1.5)}, pregnant/dramatic pauses before a loaded word, hesitant stumbles, abrupt self-interruptions or cut-offs (\texttt{---}), and forced silences (muting a panelist's mic or talking over them). & Anchor pauses \texttt{(1.5)} before saying \textit{"'loota'"} for dramatic weight; Addl. DCP cuts off mid-sentence (\textit{"wo toh--- abhi investigation chal raha hai"}) when asked if police struck first; panelist's mic faded during counter-question. & Where does the pause or stop occur? Is it deliberate/rhetorical (for suspense, emphasis, or gravitas) or defensive/evasive (cognitive load, self-repair, avoiding blame)? Who controls turn-taking and silence? \\ \hline
\textbf{2. Trailing Sentences \& Unfinished Utterances} & Clauses that trail off (\texttt{...}), drop in volume at the end, or leave the conclusion unspoken (aposiopesis) so the audience supplies the implication without the speaker taking explicit responsibility. & \textit{"Aur phir wahan jo log ghuse the... aap samajh hi sakte hain..."}; \textit{"500 salon ke sangharsh ke baad agar aisa..."} trailing off with a slow head shake. & What exact thought is left unfinished? What presupposition or shared bias is the viewer invited to fill in? Does trailing off act as a safe hedge to imply higher-level complicity or "anti-national" motive without stating it? \\ \hline
\textbf{3. Intonation, Pitch, Volume \& Prosody} & Pitch contour (rising \texttt{$\uparrow$} for incredulity/challenge; falling \texttt{$\downarrow$} for finality), volume shifts (shouting vs. hushed whisper), emphatic word stress, speech rate (rapid-fire urgency vs. slow solemnity), and tonal register (sarcastic, reverential, combative). & Hushed, reverential cadence on \textit{"Prabhu Shri Ram ke Mandir"} vs. sharp, high-pitched interrogation on \textit{"par Akhilesh ji Ram Mandir kyon$\uparrow$ nahin jate?!"}; satirical sing-song inflection on \textit{"priventiv arenjment"}. & Which exact words receive pitch peaks or volume stress? How does vocal tone turn neutral words into irony, mockery, or moral outrage? Does the speaker shift prosodic register depending on which actor is discussed? \\ \hline
\textbf{4. Coughs, Sighs \& Non-Lexical Vocalisations} & Audible sighs (\texttt{[sigh]}), scoffs/snorts (\texttt{[scoff]}), ironic chuckles (\texttt{[chuckle]}), throat-clearing/coughs (\texttt{[cough]}), sharp intakes of breath/gasps (\texttt{[gasp]}), and tongue-clicks (\texttt{[tsk]}). & Heavy \texttt{[sigh]} before reporting Trust oversight lapses (signalling reluctant sorrow rather than hostility); mocking \texttt{[scoff]} or \texttt{[chuckle]} when reading an opponent's denial; \texttt{[cough / throat-clear]} when pivoting away from an awkward fact. & What emotion or epistemic stance does the vocalisation perform (exasperation, moral disgust, disbelief, ridicule, discomfort)? Does a sigh or scoff pre-judge a speaker's claim before any verbal rebuttal is made? \\ \hline
\textbf{5. Hand Gestures, Facial Expressions \& Posture (Kinesics)} & Hand movements (air-quotes, finger-pointing/wagging, desk-thumping, open palms, counting on fingers, dismissive wave), facial cues (raised eyebrows, eye-rolls, smirk, furrowed brow), head nods/shakes, leaning into camera, and visual props (bandaged limbs, documents). & Anchor forms physical air-quotes on \textit{"'loota'"} or \textit{"'peaceful' protest"}; host points finger/pen at lens during an ultimatum; injured DCP holds bandaged hand up to camera; creator raises one eyebrow and leans in for a conspiratorial reveal. & How do gestures reinforce, contradict, or sharpen the spoken words? During split-screen listening, does the host's face (nodding vs. smirking/eye-rolling) cue the audience on whom to trust? How are bodily visuals used as evidence? \\ \hline
\end{longtable}

\subsection*{Quick Multimodal Transcription Notation Key (use inside transcript excerpts)}

\begin{longtable}{|p{0.15\textwidth}|p{0.15\textwidth}|p{0.15\textwidth}|p{0.15\textwidth}|p{0.15\textwidth}|p{0.15\textwidth}|}
\hline
Symbol & Meaning & Example in transcript & Symbol & Meaning & Example in transcript \\ \hline \endhead
\texttt{(.)} & Brief micro-pause (< 0.5 sec) & \texttt{woh (.) suparavijan lacking tha} & \texttt{$\uparrow$ / $\downarrow$} & Marked rising / falling pitch & \texttt{kyon$\uparrow$ FIR nahin karaee$\downarrow$} \\ \hline
\texttt{(1.5)} & Timed pause in seconds & \texttt{kisne (1.5) 'loota'?} & \texttt{WORD / \_word\_} & Loud volume / emphatic stress & \texttt{ye PAAP hai / \_bilkul\_ galat} \\ \hline
\texttt{---} & Abrupt stop, cut-off, or self-repair & \texttt{police ki taraf se--- abhi jaanch...} & \texttt{$^\circ$word$^\circ$} & Softened / whispered delivery & \texttt{$^\circ$ander ki baat ye hai$^\circ$} \\ \hline
\texttt{...} & Trailing off / unfinished sentence & \texttt{aap jaante hain ki kaun log...} & \texttt{[sigh] [cough] [scoff]} & Non-lexical vocalisations & \texttt{[sigh] 500 salon ke baad...} \\ \hline
\texttt{[overlap]} & Simultaneous / overlapping speech & \texttt{[Anjana: suniye---] [Guest: nahin!]} & \texttt{\{\{gesture: ...\}\}} & Hand gesture / facial cue / posture & \texttt{\{\{air-quotes on 'loota'\}\}} \\ \hline
\end{longtable}

\textbf{Methodological Reminder:} A non-verbal cue is not automatically manipulative or ideological---a cough, pause, or rapid cadence may reflect studio audio delay, live time pressure, or a speaker's habitual style. Always check (1) whether the cue clusters with verbal features, and (2) whether the speaker's delivery shifts noticeably across different actors or events (within-outlet baseline).\par\smallskip

\bigskip\hrule\bigskip

\subsection*{Event 1. Ram Mandir donation controversy (Non-Verbal \& Paralinguistic Layer)}

\begin{itemize}
  \item \textbf{Focus:} How do prosody, pauses, trailing sentences, vocalisations (sighs, coughs, scoffs), and body language frame the controversy---as solemn devotional sorrow, administrative scrutiny, or partisan combat?
  \item \textbf{Outlet:} The Wire (Senior Editorial Commentary / Analysis by Arfa Khanum Sherwani / Siddharth Varadarajan team) | \textbf{Date:} June 2021
  \item \textbf{Video title:} Ram Mandir Trust Land Deal Controversy: Serious Questions of Corruption \& Accountability | \href{https://youtu.be/7rUL8fGgRJA}{https://youtu.be/7rUL8fGgRJA} | \textbf{Timestamp(s):} \texttt{[00:00]}--\texttt{[14:28]}
  \item \textbf{Analyst:} Baani Solanki (\texttt{2025114016}) | \textbf{Language:} Hindi / Urdu-inflected Hindustani \& English (transcribed in Roman script)
  \item \textbf{Format / Visual setup:} \texttt{[ ]} TV Panel Debate \texttt{[ ]} TV Bulletin \texttt{[x]} Solo YouTube Explainer \texttt{[x]} Interview / Discussion
  \item \textbf{Primary speaker(s) observed:} Senior Journalist / Presenter at \textit{The Wire} (seated in a quiet, sober newsroom/studio setting, delivering measured editorial analysis with deed overlays and press clips)
\end{itemize}

\subsubsection*{Event-specific non-verbal questions (Event 1)}

\noindent\textbf{1.} \textbf{Vocal \& gestural marking of core labels ("scam", "loota", "paap", "gadabdi", "avyavastha"): When the controversy is named, what non-verbal cues accompany the label? (Look for physical air-quotes, vocal scare-quotes, dramatic pauses before the word, hushed/reverential drops in volume, or shouted stress.)}\par\smallskip
\begin{itemize}
  \item \textbf{Notes:}
\textit{The Wire} avoids both devotional melodrama (\texttt{"paap"}) and tabloid screaming, marking the core labels---\textit{"bhrashtachar"} (corruption), \textit{"ghotala"} (scam), and \textit{"financial impropriety"}---with \textbf{low-pitched, grave, deliberate falling intonation \texttt{$\downarrow$} and measured pauses}:\par\smallskip
  \item At \texttt{[00:45]}, the presenter pauses \texttt{(1.0)} before naming the charge: \textit{"Shri Ram Janmabhoomi Teerth Kshetra Trust ke zameen saude mein (1.0) \_gambhir\_ bhrashtachar ke aarop lage hain$\downarrow$"} (\texttt{[00:52]}), accompanied by a slow, serious head nod \texttt{\{\{gesture: slow grave nod + furrowed brow\}\}}.
  \item Conversely, when quoting the Trust's claim of \textit{"paaradarshita"} (transparency, \texttt{[06:15]}), the presenter uses \textbf{vocal scare-quotes}---raising one eyebrow, pausing \texttt{(0.5)}, and inflecting \textit{"paaradarshita$\uparrow$"} with dry, skeptical rising pitch and a subtle ironic smirk \texttt{\{\{gesture: raised eyebrow + dry half-smirk on 'paaradarshita'\}\}}.
\end{itemize}

\noindent\textbf{2.} \textbf{Differential demeanour toward actors (Temple, Trust leadership, counting clerks, donors, opposition accusers): How do the speaker's voice pitch, facial expression, and posture shift when referring to Prabhu Shri Ram / devotees vs. top figures (Champat Rai, Nripendra Misra, CM Yogi) vs. arrested counting staff vs. opposition leaders (Akhilesh Yadav, Awadhesh Prasad)?}\par\smallskip
\begin{itemize}
  \item \textbf{Notes:}
The presenter maintains a \textbf{sober, secular-institutional demeanour} across all actors, differentiating them through analytical tone rather than theatrical volume shifts:\par\smallskip
  \item \textbf{Donors / Ordinary Devotees:} Referred to with quiet civic sympathy (\texttt{"crore-on aam nagrikon aur bhakton ka paisa"}, \texttt{[01:30]}), opening both hands palm-up on the desk \texttt{\{\{gesture: both palms open resting on desk\}\}} to frame them as public stakeholders cheated by elites.
  \item \textbf{Trust Leadership \& BJP Figures (Champat Rai, Mayor Rishikesh Upadhyay, Anil Mishra):} Addressed with a \textbf{steady, unblinking prosecutorial gaze and clinical precision} (\texttt{[03:20]}--\texttt{[05:45]}). There is zero deference to their religious status; Champat Rai is treated strictly as the chief executive of a public trust created by a Supreme Court mandate (\texttt{[07:10]}).
  \item \textbf{Opposition Accusers \& Local Journalists (Sanjay Singh, Pawan Pandey, Ayodhya press):} Referred to with respectful, collegial nods (\texttt{[02:10]}), crediting them for bringing primary registration deeds into the public domain.
\end{itemize}

\noindent\textbf{3.} \textbf{Solemn devotional grief vs. high-decibel partisan spectacle: How do pauses, sighs, speech rate, and hand gestures construct the segment's emotional register? Where is delivery slow, hushed, and mournful (sighs, bowed head, folded hands), and where does it escalate into rapid-fire shouting, desk-thumping, or mocking irony?}\par\smallskip
\begin{itemize}
  \item \textbf{Notes:}
\textit{The Wire} operates in a \textbf{slow, grave, seminar-room / editorial-board register} (\textasciitilde{}115--125 words per minute) throughout the entire 14-minute video. There is zero shouting, zero desk-thumping, and zero devotional hand-folding. Instead, emotional weight is constructed through \textbf{long, reflective silences \texttt{(1.0)}--\texttt{(1.5)}, weary civic sighs \texttt{[sigh]}, and bitingly understated irony}. For example, at \texttt{[08:20]}, when discussing how ruling-party leaders invoke Lord Ram to shield themselves from audit questions, the speaker lets out a quiet \texttt{[sigh]}, shakes their head slowly, and speaks in a lowered, sorrowful register \texttt{$^\circ$...$^\circ$}: \textit{"$^\circ$Aastha ki aad mein (.) sawaal poochhne walon ko hi deshdrohi ya Ram-virodhi keh diya jaata hai$^\circ$$\downarrow$"} (\texttt{[08:28]}).\par\smallskip
\end{itemize}

\noindent\textbf{4.} \textbf{Non-verbal evaluation of financial scrutiny: When questions about unreceipted offerings, lack of frisking, or land deals are raised, how does the host/speaker react non-verbally (earnest nodding, leaning into the camera, open-palm appeals vs. eye-rolls, scoffs, dismissive hand waves, or incredulous rising pitch $\uparrow$)?}\par\smallskip
\begin{itemize}
  \item \textbf{Notes:}
Financial scrutiny is treated as a \textbf{vital constitutional and public-interest duty}. During the examination of the \texttt{18 March 2021} sale deed vs. agreement to sell (\texttt{[02:40]}--\texttt{[05:30]}), the presenter leans slightly forward toward the camera, holds up the printed registration papers / points to the screen overlay \texttt{\{\{gesture: holding printed document + tapping key clauses with pen\}\}}, and uses deliberate finger-counting (\texttt{[03:15]}) to highlight three unanswered questions: (1) the 825\% markup in 300 seconds, (2) the identical witnesses (Anil Mishra and Mayor Rishikesh Upadhyay), and (3) why the Trust did not buy directly from the original owner Kusum Pathak (\texttt{[04:50]}).\par\smallskip
\end{itemize}

\noindent\textbf{5.} \textbf{Performing certainty vs. hedging and innuendo (pauses, stops, trailing off): Are allegations delivered with definitive falling intonation $\downarrow$, finger-counting, and pointing at on-screen evidence (CCTV, documents), or hedged with hesitant pauses (.), abrupt self-interruptions (---), throat-clearing [cough], and trailing sentences (...) that imply higher-level complicity without naming it?}\par\smallskip
\begin{itemize}
  \item \textbf{Notes:}
\textit{The Wire} combines \textbf{documentary certainty} on the paper trail with \textbf{grave, trailing rhetorical reflections (\texttt{...})} on the political ecosystem:\par\smallskip
  \item \textbf{Documentary Certainty:} The dates, stamp values, and witness signatures are delivered with crisp falling intonation \texttt{$\downarrow$} and rhythmic pen-taps on the desk (\texttt{"Yeh kisi vipakshi party ka aarop nahin, balki Uttar Pradesh sarkar ke registration vibhag ke kagaz hain$\downarrow$"}, \texttt{[03:55]}).
  \item \textbf{Implicating Higher Power via Trailing Off (\texttt{...}):} At \texttt{[10:15]}, when asking how a local BJP Mayor and an RSS-nominated Trustee could execute such a transaction without top-level clearance, the presenter slows down and trails off: \textit{"Bina upar ki manzoori ke (.) Ayodhya jaise high-profile project mein itna bada khel ho jaaye... yeh baat kisi ke gale nahin utarti..."} (\texttt{[10:25]}), accompanied by a slow head shake and raised eyebrows \texttt{\{\{gesture: slow head shake + arched eyebrows\}\}}.
\end{itemize}

\noindent\textbf{6.} \textbf{Turn-taking, interruptions, and acoustic treatment of denials: In debates/interviews: who receives calm, uninterrupted speaking time and sympathetic nods, and who is cut off mid-sentence (---), talked over, or muted? In solo videos: how are clips of official denials framed via reaction pauses, raised eyebrows, or sarcastic chuckles?}\par\smallskip
\begin{itemize}
  \item \textbf{Notes:}
In this editorial commentary format, there is no live shouting match. However, when Champat Rai's brief video clip dismissing the allegations (\texttt{"hum is par koi tippani nahin karenge"}, \texttt{[07:35]}) is played, the presenter follows it with a \textbf{pregnant \texttt{(1.5)} second silence}, looking directly into the camera with a raised eyebrow and a quiet, dry \texttt{[scoff]} (\texttt{[07:48]}): \textit{"Jawab dene ke bajaye (1.5) maun ya phir dhamki$\downarrow$"} (\texttt{[07:52]}). The silence frames the Trust's refusal to answer questions as an admission of guilt.\par\smallskip
\end{itemize}

\noindent\textbf{7.} \textbf{Performing the boundary between religious devotion and institutional critique: How does the presenter use vocal tone, sighs, and body language to signal personal reverence for the Temple (or secular objectivity) while criticising---or defending---the Trust's financial management?}\par\smallskip
\begin{itemize}
  \item \textbf{Notes:}
Unlike Aaj Tak or Republic TV (which perform Hindu devotional insiderhood), \textit{The Wire} explicitly performs \textbf{secular, legal-constitutional objectivity}. The presenter sits with composed, journalistic posture, keeps hands folded calmly on the desk between points \texttt{\{\{gesture: hands lightly clasped on desk in newsroom posture\}\}}, and grounds the critique in legal accountability: because the Shri Ram Janmabhoomi Teerth Kshetra Trust was constituted under a \textbf{Supreme Court judgment} (\texttt{[07:10]}) and receives tax-exempt public donations, its finances are a matter of secular public law, not private religious faith (\texttt{[07:25]}).\par\smallskip
\end{itemize}

\noindent\textbf{8.} \textbf{Self-presentation as the "guardian of public trust": How does the speaker use direct eye-gaze, posture, moral sighs, or emphatic gestures (pointing at the viewer, hands on heart/desk) to position themselves as the authentic protector of devotees or public accountability?}\par\smallskip
\begin{itemize}
  \item \textbf{Notes:}
The presenter positions \textit{The Wire} as the \textbf{uncompromising fourth estate ("speaking truth to power")} at a time when mainstream television has surrendered its watchdog role. Using steady, unhurried eye-gaze into the lens, a weary moral \texttt{[sigh]} over attacks on local Ayodhya journalists who reported the scam (\texttt{[11:40]}), and a calm open-palm appeal (\texttt{[13:50]}), the host projects quiet moral authority rather than theatrical rage.\par\smallskip
\end{itemize}

\subsubsection*{Apply the five non-verbal features (Event 1)}

\begin{longtable}{|p{0.31\textwidth}|p{0.31\textwidth}|p{0.31\textwidth}|}
\hline
Non-verbal feature & Timestamp + spoken anchor words + non-verbal notation & What it does in context (stance, framing, implication) \\ \hline \endhead
\textbf{1. Pauses, silences \& stops} & \texttt{[07:48]} \texttt{[(1.5) silent reaction pause after Champat Rai clip]} \textit{"Jawab dene ke bajaye (1.5) maun ya phir dhamki$\downarrow$"} & The \texttt{(1.5)} silence lets Champat Rai's evasive response hang in the air, framing his refusal to answer as arrogant stonewalling. \\ \hline
\textbf{2. Trailing sentences (\texttt{...})} & \texttt{[10:25]} \textit{"Bina upar ki manzoori ke Ayodhya jaise project mein itna bada khel ho jaaye... yeh baat kisi ke gale nahin utarti..."} & Leaves the identity of \textit{"upar"} (top state/national leadership) unspoken while inviting the viewer to infer high-level political sanction. \\ \hline
\textbf{3. Intonation, pitch \& volume} & \texttt{[00:52]} \textit{"\_Gambhir\_ bhrashtachar ke aarop lage hain$\downarrow$"} (low, grave falling pitch \texttt{$\downarrow$}) vs. \texttt{[06:15]} \textit{"Trust ki 'paaradarshita'$\uparrow$"} (skeptical rising inflection \texttt{$\uparrow$}) & Grave falling intonation lends institutional weight to the corruption charge, while rising inflection on \texttt{'paaradarshita'} turns the Trust's claim into irony. \\ \hline
\textbf{4. Coughs, sighs \& vocalisations} & \texttt{[07:49]} \texttt{[[scoff]]} (soft dry scoff after Trust denial) / \texttt{[08:20]} \texttt{[[sigh]]} \textit{"$^\circ$Aastha ki aad mein sawaal poochhne walon ko hi...$^\circ$"} & The soft \texttt{[scoff]} punctures the Trust's moral posturing; the weary \texttt{[sigh]} performs journalistic sorrow over the weaponisation of faith. \\ \hline
\textbf{5. Hand gestures, face \& posture} & \texttt{[03:15]} \texttt{\{\{gesture: holding printed registry papers + tapping 3 points with pen\}\}} / \texttt{[06:15]} \texttt{\{\{gesture: single raised eyebrow + dry half-smirk\}\}} & Holding physical papers grounds the critique in public records; the raised eyebrow cues skeptical distance from official statements. \\ \hline
\end{longtable}

\subsubsection*{Short summary (Event 1 Non-Verbal Synthesis)}

\begin{itemize}
  \item \textbf{What overall interpretation or emotional stance does the non-verbal delivery encourage?}
The non-verbal delivery encourages \textbf{sober, critical, secular-constitutional scrutiny}. By eschewing both devotional emotionalism and studio shouting in favour of a slow (\textasciitilde{}120 WPM), low-decibel newsroom cadence, long \texttt{(1.5)} reaction silences, and methodical document-tapping, \textit{The Wire} frames the land transaction as a serious abuse of public trust shielded by religious majoritarianism.\par\smallskip
  \item \textbf{Whose voice or physical presence is acoustically/visually foregrounded, undermined, or silenced?}
\textbf{The registration deeds, local Ayodhya investigative journalists, and opposition whistleblowers} are visually and acoustically foregrounded with calm authority. \textbf{Champat Rai and the Trust leadership} are acoustically and visually undermined---their brief denial clips are followed by pregnant \texttt{(1.5)} silences, raised eyebrows, and dry scoffs that expose their refusal to engage with the documentary facts.\par\smallskip
  \item \textbf{What evidence supports your non-verbal reading? What complicates it (e.g., speaker baseline, format)?}
The contrast between the crisp falling intonation \texttt{$\downarrow$} and pen-tapping on registration facts (\texttt{[03:15]}--\texttt{[03:55]}) versus the rising skeptical inflection \texttt{$\uparrow$} on \textit{"paaradarshita"} (\texttt{[06:15]}) and trailing off \texttt{...} on higher political approval (\texttt{[10:25]}) shows how subtle prosody replaces loud volume in serious editorial commentary. A complicating factor is \textbf{speaker baseline}: \textit{The Wire}'s presenters consistently use this grave, understated, slightly mournful cadence across all political broadcasts, meaning the low-key delivery is an institutional house style designed to signal "serious journalism" in contrast to mainstream TV.\par\smallskip
\end{itemize}

\bigskip\hrule\bigskip

\subsection*{Event 2. Students' protest regarding NEET paper leak (Non-Verbal \& Paralinguistic Layer)}

\begin{itemize}
  \item \textbf{Focus:} How do prosody, pauses, trailing sentences, vocalisations, and visual body language construct the legitimacy of student protesters, police use of force, and state accountability?
  \item \textbf{Outlet:} The Wire (Ground Report \& Studio Analysis on NEET Paper Leak \& Student Protests) | \textbf{Date:} June 2024
  \item \textbf{Video title:} NEET Paper Leak Protest: Students March, Police Crackdown \& NTA Accountability | \href{https://youtu.be/pYaTVvHeGoM}{https://youtu.be/pYaTVvHeGoM} | \textbf{Timestamp(s):} \texttt{[00:00]}--\texttt{[16:50]}
  \item \textbf{Analyst:} Baani Solanki (\texttt{2025114016}) | \textbf{Language:} Hindi \& English (transcribed in Roman script)
  \item \textbf{Format / Visual setup:} \texttt{[ ]} TV Panel Debate \texttt{[ ]} TV Bulletin \texttt{[ ]} Solo YouTube Explainer \texttt{[x]} Interview / Discussion
  \item \textbf{Primary speaker(s) observed:} \textit{The Wire} Ground Reporter \& Studio Anchor, alongside prolonged, uninterrupted on-street interviews with protesting NEET Students, Youth Activists, and Parents at Jantar Mantar / Parliament Street
\end{itemize}

\subsubsection*{Event-specific non-verbal questions (Event 2)}

\noindent\textbf{1.} \textbf{Prosodic \& gestural marking of participant labels ("students/bachche" vs. "protesters" vs. "upadravi/goons"): How do vocal tone, volume, facial expression, and hand gestures shift across labels? (e.g., softened, empathetic pitch for "students" vs. harsh, staccato stress, scowls, or accusatory finger-pointing on "upadraviyon", "goons", or "anti-social elements".)}\par\smallskip
\begin{itemize}
  \item \textbf{Notes:}
\textit{The Wire} never uses state-policing labels like \textit{"upadravi"} or \textit{"goons"}. Protesters are consistently referred to as \textit{"chhatra"} (students), \textit{"aabhyarthi"} (aspirants), \textit{"naujawan"} (youth), and \textit{"satyagrahi / pradarshankari"} (\texttt{[00:30]}, \texttt{[04:15]}, \texttt{[11:20]}) with a \textbf{warm, respectful, level vocal tone and open, inviting hand gestures} holding the microphone toward them \texttt{\{\{gesture: reporter extending mic steadily with empathetic nod\}\}}. When quoting ruling-party or media claims that protesters are \textit{"political elements"} (\texttt{[12:40]}), the studio anchor marks the phrase with \textbf{skeptical rising intonation \texttt{$\uparrow$} and raised eyebrows}, visually rejecting the criminalisation of student dissent.\par\smallskip
\end{itemize}

\noindent\textbf{2.} \textbf{Delivery of students' stated demands: When the protesters' core demands (Education Minister's resignation, compensation, release of Sonam Wangchuk, exam reform) are mentioned, are they delivered slowly with deliberate pauses and finger-counting, rushed through in a flat monotone, mocked with air-quotes/scoffs, or trailed off (...)?}\par\smallskip
\begin{itemize}
  \item \textbf{Notes:}
The protesters' core demands---(1) immediate resignation of Education Minister Dharmendra Pradhan, (2) abolition/complete restructuring of the NTA, (3) re-NEET / fair justice for affected candidates, and (4) stopping police repression of peaceful activists (\texttt{[03:10]}--\texttt{[05:45]}, \texttt{[13:15]})---are delivered \textbf{both directly by the students in their own voices and summarised slowly by the host with deliberate \texttt{(0.8)} pauses and finger-counting} \texttt{\{\{gesture: counting demands on fingers with firm nods\}\}}. Far from mocking the demand for the Education Minister's resignation, the host delivers it with \textbf{grave falling intonation \texttt{$\downarrow$}} (\textit{"Chhatron ka seedha sawaal hai---Dharmendra Pradhan istifa kyon nahin dete$\downarrow$"}, \texttt{[13:28]}), treating ministerial resignation as standard democratic accountability.\par\smallskip
\end{itemize}

\noindent\textbf{3.} \textbf{Where does vocal \& kinetic energy peak---grievance, street disruption, or electoral horse-race?: Track where the speaker's volume, pitch range, pace, and gesture intensity reach their climax: during accounts of the paper leak/student suicides, during street clashes and barricade breaches, or during partisan BJP-vs-Opposition sparring.}\par\smallskip
\begin{itemize}
  \item \textbf{Notes:}
Vocal and kinetic energy peaks during \textbf{the students' own raw, emotional testimonies on the street (\texttt{[02:15]}--\texttt{[06:30]}) and during ground footage of police dragging and detaining peaceful protesters (\texttt{[07:40]}--\texttt{[10:15]})}:\par\smallskip
  \item When a female NEET aspirant at the protest site speaks with a trembling voice, rising pitch \texttt{$\uparrow$}, and tears in her eyes (\textit{"Humne \texttt{3} saal apni zindagi ke diye hain, aur badle mein humein lathiyan mil rahi hain$\uparrow$!"}, \texttt{[03:42]}), the camera holds a steady, unflinching eye-level close-up while the reporter listens in complete silence with a solemn, empathetic nod.
\end{itemize}

\noindent\textbf{4.} \textbf{Non-verbal attribution of responsibility (sighs, stops, head shakes, pointing): When examination authorities (NTA), the Education Minister, police, opposition politicians, or students are named as responsible, what non-verbal cues accompany the attribution (exasperated sighs, mocking laughs, abrupt self-stops ---, or direct finger-pointing at the lens)?}\par\smallskip
\begin{itemize}
  \item \textbf{Notes:}
Responsibility is attributed directly to the \textbf{Union Education Minister (Dharmendra Pradhan), the NTA leadership (Subodh Kumar Singh), and the ruling establishment's pattern of institutional capture}:\par\smallskip
  \item When discussing how the NTA chief and Education Minister initially denied any paper leak despite the Patna FIR (\texttt{[11:10]}), the studio anchor lets out an audible, weary \texttt{[sigh]}, shakes their head slowly \texttt{\{\{gesture: slow head shake + furrowed brow\}\}}, and taps the desk with a pen on each contradiction (\texttt{[11:25]}).
  \item When referencing the Godhra (Gujarat) and Hazaribagh/Patna exam centre nexus (\texttt{[12:05]}), the host raises an eyebrow and uses pointed downward hand gestures to highlight the political connections of the centre operators.
\end{itemize}

\noindent\textbf{5.} \textbf{Prosody of questions---open inquiry vs. loaded rhetorical challenge: Are questions about the protest delivered with level, open intonation, or with sharp rising pitch $\uparrow$, sarcastic stress, raised eyebrows, and trailing clauses (...) that presuppose protesters are violent or politically manipulated?}\par\smallskip
\begin{itemize}
  \item \textbf{Notes:}
\textit{The Wire}'s prosody of questions is the \textbf{exact mirror image of Aaj Tak}:\par\smallskip
  \item \textbf{Questions to Protesting Students (\texttt{[02:10]}, \texttt{[04:50]}):} Delivered with \textbf{softened, level, open inquiry intonation} (\texttt{"Aap log itni garmi mein yahan baithe hain, aapki mukhya maangein kya hain$\downarrow$?"}, \texttt{[02:14]}), giving the student full agency to frame the issue.
  \item \textbf{Questions Directed at the State / Police / NTA (\texttt{[08:50]}, \texttt{[13:25]}):} Delivered with \textbf{pointed rising pitch \texttt{$\uparrow$}, emphatic stress, and prosecutorial gravity} (\texttt{"Kya apna haq maangne wale chhatron par lathi chalana aur unhe hirasat mein lena hi sarkar ka jawab hai$\uparrow$?"}, \texttt{[08:55]}).
\end{itemize}

\noindent\textbf{6.} \textbf{Visual \& paralinguistic performance of injury and violence: How are claims of violence substantiated non-verbally on screen? (Compare camera close-ups and pained vocal tone on police bandages/fractured hands vs. visual/acoustic treatment of student head injuries, lathi-charge footage, or defensive stammers/pauses when an officer or guest is asked who used force first.)}\par\smallskip
\begin{itemize}
  \item \textbf{Notes:}
Whereas Aaj Tak substantiates violence by zooming in on a police officer's bandaged hand, \textit{The Wire} substantiates violence through \textbf{eye-level camera close-ups of bruised students, torn clothes, and police personnel grabbing protesters by the limbs and shoving them into buses} (\texttt{[07:45]}--\texttt{[09:30]}). A student protester holds up his scraped arm/wrist to \textit{The Wire}'s camera \texttt{\{\{gesture: student holds bruised wrist up to camera lens\}\}} (\texttt{[08:12]}) while breathing heavily \texttt{[gasp]} and saying: \textit{"Dekhiye police ne kaise ghasita hai$\downarrow$"} (\texttt{[08:15]}). The reporter lowers their voice in solidarity (\texttt{$^\circ$aapko chot aayi hai$^\circ$}, \texttt{[08:18]}), visually and acoustically validating student injury at the hands of state force.\par\smallskip
\end{itemize}

\noindent\textbf{7.} \textbf{Asymmetric delivery of police actions vs. protester actions: Compare the vocal cadence, pauses, and hand gestures used when describing police actions (e.g., calm, matter-of-fact administrative tone on "tear gas chhode gae" / "preventive arrangement") versus protester actions (alarmist pitch, striking gestures, or ironic softening as "retaliation").}\par\smallskip
\begin{itemize}
  \item \textbf{Notes:}
  \item \textbf{Police Actions:} Described in an \textbf{indignant, critical cadence} that strips away bureaucratic euphemisms: when quoting the police claim of \textit{"preventive detention"} or \textit{"law and order arrangement"} (\texttt{[09:10]}), the presenter uses a \textbf{satirical, dry inflection and physical air-quotes} \texttt{\{\{gesture: air-quotes on 'preventive detention'\}\}}, contrasting the euphemism with the visual footage of students being dragged across the asphalt.
  \item \textbf{Protester Actions:} Described in a \textbf{dignified, empathetic cadence} (\texttt{"shantipurna tarike se Sansad gherao karne ja rahe chhatra"}, \texttt{[07:50]}), framing their march as a constitutional right to peaceful assembly.
\end{itemize}

\noindent\textbf{8.} \textbf{Acoustic \& visual space granted to students' own voices: Are protesting students heard in their own voice with clear audio, steady close-ups, and empathetic listening nods, or are they reduced to noisy background chaos B-roll while studio anchors/politicians speak over them?}\par\smallskip
\begin{itemize}
  \item \textbf{Notes:}
Protesting students receive \textbf{over 60\% of the total speaking time in the segment} with \textbf{clear, unedited close-up audio and zero anchor \texttt{[overlap]} interruptions}. During each student interview (\texttt{[02:15]}--\texttt{[06:30]}, \texttt{[08:05]}--\texttt{[10:00]}), the reporter holds the microphone steadily, maintains eye contact, nods empathetically \texttt{\{\{gesture: continuous empathetic listening nods\}\}}, and waits a full \texttt{(0.8)}--\texttt{(1.0)} second after the student finishes speaking before asking a follow-up question---the exact opposite of television panel cross-talk.\par\smallskip
\end{itemize}

\noindent\textbf{9.} \textbf{Non-verbal shift between the NEET paper-leak issue and the street protest: Does the speaker's non-verbal demeanour shift from sympathetic concern (sombre tone, sighs, solidarity gestures) when discussing the exam leak itself to combative, authoritarian, or dismissive delivery once students march toward Parliament?}\par\smallskip
\begin{itemize}
  \item \textbf{Notes:}
\textbf{No authoritarian shift occurs toward the students.} Instead, as the video moves from the exam paper leak (\texttt{[00:00]}--\texttt{[02:00]}) to the street march and police crackdown (\texttt{[07:40]}--\texttt{[10:30]}), the presenter's non-verbal demeanour shifts from \textbf{sombre sympathy for exam victims to grave constitutional outrage against police high-handedness}, treating the street protest as the natural democratic extension of the exam grievance.\par\smallskip
\end{itemize}

\subsubsection*{Apply the five non-verbal features (Event 2)}

\begin{longtable}{|p{0.31\textwidth}|p{0.31\textwidth}|p{0.31\textwidth}|}
\hline
Non-verbal feature & Timestamp + spoken anchor words + non-verbal notation & What it does in context (stance, framing, implication) \\ \hline \endhead
\textbf{1. Pauses, silences \& stops} & \texttt{[03:48]} \texttt{[(1.0) respectful silence by reporter after crying student's testimony]} / \texttt{[11:20]} \textit{"Pehle kaha koi leak nahin hua--- (1.2) aur ab CBI jaanch$\downarrow$"} & The \texttt{(1.0)} silence honours the student's raw emotion without studio intrusion; the \texttt{(1.2)} pause highlights the government's U-turn on the paper leak. \\ \hline
\textbf{2. Trailing sentences (\texttt{...})} & \texttt{[14:15]} \textit{"Jab har badi pariksha ka paper leak ho raha ho, aur sawaal poochhne par chhatron ko lathiyan milein... toh is desh ke yuvaon ka bhavishya..."} & Trails off with a slow, sorrowful head shake, inviting the viewer to reflect on authoritarian intolerance toward youth dissent. \\ \hline
\textbf{3. Intonation, pitch \& volume} & \texttt{[02:14]} \textit{"$^\circ$Aapki mukhya maangein kya hain$^\circ$$\downarrow$?"} (gentle, open inquiry to student) vs. \texttt{[08:55]} \textit{"Kya chhatron par \_LATHI\_ chalana hi sarkar ka jawab hai$\uparrow$?!"} & Gentle, low-volume open prosody empowers student speakers, whereas sharp rising pitch \texttt{$\uparrow$} and stress on \texttt{\_LATHI\_} indict state repression. \\ \hline
\textbf{4. Coughs, sighs \& vocalisations} & \texttt{[08:12]} Student: \texttt{[[gasp]]} \textit{"Dekhiye police ne kaise ghasita hai$\downarrow$"} / \texttt{[11:10]} Anchor: \texttt{[[sigh]]} (on Education Minister's initial denial) & Including the student's breathless \texttt{[gasp]} on camera conveys authentic physical distress; the anchor's \texttt{[sigh]} performs civic exhaustion with official denialism. \\ \hline
\textbf{5. Hand gestures, face \& posture} & \texttt{[08:12]} \texttt{\{\{gesture: student holds bruised wrist up to camera lens\}\}} / \texttt{[09:10]} \texttt{\{\{gesture: anchor air-quotes on 'preventive detention'\}\}} / \texttt{[02:20]} \texttt{\{\{gesture: reporter's steady mic hold + empathetic nods\}\}} & The student's bruised wrist directly counters TV visuals of police bandages; air-quotes debunk police euphemisms; empathetic nods validate student testimony. \\ \hline
\end{longtable}

\subsubsection*{Short summary (Event 2 Non-Verbal Synthesis)}

\begin{itemize}
  \item \textbf{What overall interpretation or emotional stance does the non-verbal delivery encourage?}
The non-verbal delivery encourages \textbf{solidarity with protesting students and moral indignation at both the NTA's exam failure and the police crackdown on peaceful marchers}. Through eye-level close-ups, unmuted student audio, and grave studio commentary, the viewer is guided to see the protesters as courageous citizens defending public education against an unaccountable state.\par\smallskip
  \item \textbf{Whose voice or physical presence is acoustically/visually foregrounded, undermined, or silenced?}
\textbf{Protesting NEET aspirants and youth activists} are acoustically and visually foregrounded, receiving uninterrupted close-ups, empathetic nods, and visual validation of their bruises (\texttt{[08:12]}). \textbf{Police claims of "preventive arrangement" and the Education Minister's denials} are undermined via air-quotes, weary \texttt{[sigh]} vocalisations, and direct juxtaposition with crackdown footage.\par\smallskip
  \item \textbf{What evidence supports your non-verbal reading? What complicates it (e.g., speaker baseline, format)?}
The zero-interruption turn-taking structure during street interviews (\texttt{[02:15]}--\texttt{[06:30]}) and the visual close-up on a student's scraped wrist (\texttt{[08:12]}) provide concrete multimodal evidence of a citizen-centred frame. What complicates a purely neutral reading is that \textit{The Wire}'s ground-report format selectively foregrounds articulate, injured protesters while granting minimal screen time to police spokespersons, reflecting its counter-establishment editorial baseline.\par\smallskip
  \item \textbf{How do the non-verbal cues interact with the spoken words? (Do they reinforce, dramatise, soften/hedge, invert meaning via sarcasm/irony, or leave an accusation unspoken via trailing off?)}
The non-verbal cues \textbf{reinforce student testimonies} through silent, respectful listening nods and \texttt{(1.0)} pauses, \textbf{invert state policing euphemisms via irony} (air-quotes and skeptical pitch on \textit{"'preventive detention'"} at \texttt{[09:10]}), and \textbf{leave a broader indictment of democratic backsliding unspoken via trailing off} (\texttt{...} at \texttt{[14:15]}).\par\smallskip
  \item \textbf{What context, speaker baseline, or format factors support or complicate this interpretation? (e.g., Is this the speaker's normal style across all videos, or a marked shift on this specific actor/claim?)}
Across both Event 1 and Event 2, \textit{The Wire} maintains a consistent \textbf{low-decibel, grave, civil-liberties baseline} (\textasciitilde{}120 WPM, minimal shouting, heavy reliance on pauses \texttt{(1.5)}, weary \texttt{[sigh]} vocalisations, and skeptical eyebrow raises). However, Event 2 incorporates high-arousal ground footage of breathless \texttt{[gasp]}, crying students (\texttt{[03:42]}, \texttt{[08:12]}), creating a powerful multimodal contrast between raw street distress and calm studio indictment.\par\smallskip
  \item \textbf{How does another outlet perform the same moment, actor, or claim non-verbally?}
Compare \textbf{The Wire} (\texttt{Event 2}, \texttt{[08:12]}) with \textbf{Aaj Tak} (\texttt{Event 2}, \texttt{[08:40]}): at the exact same protest, Aaj Tak puts an Additional DCP on screen holding up his \textbf{bandaged hand} while the anchor lowers her voice in sympathy (\texttt{$^\circ$DCP sahab ke haath mein fracture hai$^\circ$}) and shouts over students \texttt{[overlap]} as \textit{"upadravi"}. \textit{The Wire} shows the exact inverse visual prop---a \textbf{student holding up his bruised wrist} after being dragged into a police bus (\texttt{[08:12]})---while the anchor uses air-quotes to mock police claims of \textit{"preventive arrangement"}.\par\smallskip
  \item \textbf{My interpretation, with reasons:}
\textit{The Wire}'s multimodal discourse across Events 1 and 2 exemplifies \textbf{counter-hegemonic journalistic framing}. In Event 1, it uses secular newsroom posture, document-tapping, and pregnant \texttt{(1.5)} silences to strip the Ram Mandir Trust of its religious immunity and demand a public audit. In Event 2, it uses acoustic generosity (uninterrupted mic time, empathetic nods) and visual counter-evidence (student bruises vs. police bandages) to legitimise street protest as an essential democratic right.\par\smallskip
\end{itemize}

\end{document}
